% Pursuing Stacks (cotes 134-2 to 134-8): archive cote 134-4, site page 38 (PDF page 39), typescript page 126.
% typescript. Draft compared: ps3.tex:1891-1967.
% AI-transcribed by Claude Opus 5.5 (claude-opus-5-5, adaptive thinking, effort high), Message Batches API, batch ps2 (msgbatch_01NVwr1Yuo8qJf5rykeGBXPK), 2026-11-12; image review recorded in records/CHECK_DECISIONS_134_3_4.tsv.
\scanpage{134-4/38}
\mspage{126}
a)\quad T\over{H}{h}e functor M\hand{$\to$}M' is bijective on objects,  surjective on arrows.
b)\quad The functor M\hand{$\to$}M' is a localization functor \hand{(}i\over{t}{t} w\over{ill}{ill} be so then,  in
view of a),  in the strict sense,  namely M' will be \inserted[typed]{M-}isomorphic,  no\hand{t}
only \inserted[typed]{M-}equivalent, to a localization $W^{-1}M$ ).

However,  the question arises w\over{\unsure{h}}{h}\over{e}{e}ther it is possible to define \over{s}{s}uch
a homotopy structure on M,  in terms of an arbitrary R i.e. an arbitrary
quotient category having same objects  (let Q(M)=Q the set of all such R's)\hand{,}
or in terms of an arbitrary localization of M,  or what amounts to the sam
% NOTE: the right edge of the scan cuts the ends of some lines (``sam'', ``definin'', ``relatio'', ``clear'', ``injectiv'')
me,  in terms of an arbitrary strongly saturated W\hand{$\subset$}$F\over{l}{\ell}$(M)  (let's call
L(M)=L  the set of all such W's).  T$\overset{\hand{\scriptscriptstyle\perp\perp}}{\underset{\hand{\scriptscriptstyle,\frown}}{\mathrm h}}$e first thing that comes to mind
here,  is that we got two natural maps
% Hand-drawn arrows are marked blue; typed objects remain plain.
\hand{(1)} \diagram{\begin{tikzcd} Q \arrow[r, bend left, "r", color=blue] \& L \arrow[l, bend left, "s", color=blue] \end{tikzcd}}
between Q and L,  which are defined by the observation that whenever we
have a functor \inserted{i:} M\hand{$\to$}M' , \inserted[typed]{injective on ob\over{j}{j}ects\hand{,}} it defines both an R\hand{$\in$}Q  (namely f\hand{$\sim$}$_{R}$g  iff
i(f)=i(g))  and a W\hand{$\in$}L (namely f\hand{$\in$}W  iff u(f) is invertible).  \over{\unsure{f}}{F}or defining
r(R)  resp. s(\over{\unsure{W}}{W}),  we apply this to the case when M' = M/R  resp. $W^{-1}$M.
Let's look a little at the two cases separately.

Start with R \inserted[typed]{in Q}, we get W=r(R),
\[ W = \hand{\{} f \hand{\in} F\over{l}{\ell}(M) \hand{\mid} i(f) \text{ invertible}\hand{\}} \hand{,} \]
where
\[ i:\ M \hand{\to} M/R \]
is the canonical functor;  we thus get a canonical functor (compatible
with the structures ``under M'' )
\[ (\over{*}{2}) \qquad \deleted{\xout{\illegible{1}}}\hand{\ensuremath{\alpha=\alpha_{R}}}:\ M_W \hand{\to} M_R \hand{,} \]
where for simplicity I write
\[ M_R = M/R \ ,\ \ M_W = W^{-1}M \ . \]
We may define W as the largest element in L (for the natural order relatio\unsure{n}
in L namely inclusion of subsets of $F\over{l}{\ell}$(M))  such that a functor (*)
exists (compatible with the functors from M into both sides)  -  such a
functor of course is unique (by the preliminaries on functors M\hand{$\to$}M' made
at the beginning).  \xout{Now} In terms of (\over{\unsure{*}}{2}),  we can say that R is \inserted{actually} a homotopy
relation  (let's call $Q_o$(M)=$Q_o$ the \xout{corresponding} subset of Q\deleted{\unsure{'s}} of all such
\xout{Now start with}
relations)  iff \deleted{(*} (\hand{2}) is an isomorphism, or equivalently (as it is clear
ly bijective on \xout{arrows} objects,  surjective on arrows)  if\hand{f} it is \uline{injectiv\hand{e}
on arrows}\hand{,} i.e. \uline{faithful}.

Conversely,  start now with W in L ,  we get R=s(W) ,
\[ R = \hand{\{}(f,g)\hand{\in} F\over{l}{\ell}(M)\mathrm{x}F\over{l}{\ell}(M) \hand{\mid} i(f)=i(g)\hand{\}} \ , \]
\margin{38}
